\subsection{A closed finite definition}

The following construction fixes all four integers in \eqref{eq:ports}.
There is no free program index, input, or oracle. Its finite local
predicate is the literal K215 predicate of Section~\ref{sec:local};
the complete transition table and predicate are included in the archive.
The coefficient definitions terminate after bounded finite enumeration,
although their binary expansions are infeasibly large.

Set \(a_0=439\) and \(\kappa=38\,256\,607\,897\).
Let \(W_0=S\), \(W_1=E_*\), with
\[
 S=(438,434,435,436,1),\qquad
 E_*=(438,0,437,438,437).
\]
List the remaining valid crosses in lexicographic order on
\(\{0,\ldots,438\}^5\), excluding these first two.
Coordinates are always prior, left, center, right, following.
The count \(\kappa\) is the exact count before the old compiler's
free-corner factor \(439^4\). The source census and the independent
replay are included; no halting test is performed in this enumeration.

Let \(\ell\) be the least integer at least two satisfying
\(2^\ell>\kappa+1\), and put \(A_0=2^\ell\). Define
\begin{align}
 p_0&=\max\left\{0,
       \left\lceil\frac{\kappa-2\ell-9a_0-4}{2}\right\rceil\right\},
 &j_{\max}&=5a_0+4+p_0,\notag\\
 \mu&=A_0-2+\sum_{j=1}^{j_{\max}}A_0^j,
 &m_0&=2\ell+14a_0+8+2p_0,\notag\\
 n_0&=m_0+1,
 &n_{\rm dum}&=n_0-\kappa-5a_0-4,\notag\\
 e_*&=\kappa+12a_0,
 &M_0&=\kappa+15a_0+n_{\rm dum}.
 \label{eq:cross-layout}
\end{align}
In this instance \(n_{\rm dum}=1\). The definitions retain the
original geometric separation between selectors, copies, dummies,
and anchors.

Use the ordered shape
\[
 \Sigma=((0,1),(1,0),(1,1),(1,2),(2,1)).
\]
For \((r,c)=\Sigma_z\), \(0\le z\le4\), and \(0\le u<a_0\), set
\begin{equation}\label{eq:cross-copy-address}
 e(r,c,u)=\kappa+3a_0+(2-r)3a_0+(2-c)a_0+u,
 \qquad j(z,u)=1+za_0+u.
\end{equation}
The allowed radix positions form the disjoint set
\begin{align}
 \mathcal E={}&\{0,\ldots,\kappa-1\}
 \ \cup\ \{e(r,c,u):(r,c)\in\Sigma,\ 0\le u<a_0\}\notag\\
 &{}\cup\{e_*,\ldots,e_*+n_{\rm dum}-1\}
 \ \cup\{M_0,3M_0,9M_0,27M_0\}.
 \label{eq:cross-support}
\end{align}
It has cardinality \(n_0\). The five copy blocks, in the order
\(\Sigma\), start at
\(\kappa+10a_0,\kappa+8a_0,\kappa+7a_0,\kappa+6a_0,\kappa+4a_0\);
each has length \(a_0\). All nonanchor positions lie below \(M_0\).

Define clause coefficients \(c_e\), for \(e\in\mathcal E\), by
\begin{align}
 c_i&=1+\sum_{z=0}^4 A_0^{\,1+za_0+W_i[z]}
            +\sum_{t=1}^4 A_0^{5a_0+t}
       &&(0\le i<\kappa),\notag\\
 c_{e(r,c,u)}&=A_0^{\,j(z,u)}
       &&((r,c)=\Sigma_z),\notag\\
 c_{3^{t-1}M_0}&=A_0^{5a_0+t}
       &&(1\le t\le4),\notag\\
 c_e&=0&&\text{at dummy positions}.
 \label{eq:cross-clause-coefficients}
\end{align}
The leading \(1\) counts selected crosses; the other digits pair each
selected cross with its five copies and four anchors.

\subsection{Band separation, radix, and the four ports}

Put
\begin{align}
 E_{\max}&=27M_0,&H_C&=51M_0+3a_0+1,\notag\\
 T_1&=H_C+2E_{\max}+a_0+1,
 &T_2&=T_1+4E_{\max}+2.
 \label{eq:cross-bands}
\end{align}
Let \(L_0\) be the least power of five at least
\(T_2+2E_{\max}+2\). For \(v>0\), let
\(\operatorname{bits}(v)=\lfloor\log_2v\rfloor+1\), and set
\begin{align}
 U_0&=4(n_0+2)\bigl(28(\kappa+5a_0+4)+6\bigr)+8,\notag\\
 b_{\min}&=\max\bigl\{
       \ell(j_{\max}+1)+3,\
       \ell(5a_0+4)+\operatorname{bits}(U_0)+1,\
       4\bigr\}.
 \label{eq:cross-radix-bound}
\end{align}
Let \(b\) be the least power of five at least \(b_{\min}\), and put
\[
 \rho=2^b,\qquad d=bL_0,\qquad B=\rho^{L_0}=2^d.
\]
All exponents used below are nonnegative.

The transport coefficients are
\begin{align}
 D_C={}&\rho^{3a_0}+\rho^{H_C+a_0}
       +\rho^{8M_0}+\rho^{24M_0}
       +\sum_{e\in\mathcal E}c_e
              \bigl(\rho^{T_1-e}+\rho^{T_2-e}\bigr),\notag\\
 D_R={}&\rho^{H_C}.
 \label{eq:cross-transport-constants}
\end{align}
Define the unshifted field mask by
\begin{align}
 M_F^0={}&\mu(\rho^{T_1}+\rho^{T_2})
       +\sum_{u=0}^{a_0-1}
            \bigl(\rho^{e(0,1,u)}+\rho^{e(1,1,u)}\bigr)\notag\\
 &+\sum_{u=0}^{a_0-1}
            \bigl(\rho^{H_C+e(1,0,u)}+\rho^{H_C+e(1,1,u)}\bigr)
       +\rho^{9M_0}+\rho^{27M_0}.
 \label{eq:cross-field-mask}
\end{align}
Finally set
\begin{equation}\label{eq:cross-word-mask}
 M_C=B-1-\sum_{\substack{e\in\mathcal E\\e\ne1}}\rho^e
                  -2\rho^{e_*},
 \qquad M_F=M_F^0+4.
\end{equation}
The four fixed polynomial coefficients are exactly
\begin{equation}\boxed{
 \beta=B-1,\qquad K=D_C+BD_R,\qquad M_C,\qquad
 \MF=M_F+B-1.}\label{eq:new-four-ports}
\end{equation}
In particular, the source port \(\MF\) is not the native mask \(M_F\).
Every arithmetic presentation in this article uses this same quartet.

The concrete scalar geometry is
\begin{center}\small
\begin{tabular}{@{}lr@{\qquad}lr@{}}
\toprule
Quantity&Value&Quantity&Value\\
\midrule
\(\ell\)&36&\(n_{\rm dum}\)&1\\
\(p_0\)&19\,128\,301\,935&\(j_{\max}\)&19\,128\,304\,134\\
\(m_0\)&38\,256\,610\,096&\(M_0\)&38\,256\,614\,483\\
\(e_*\)&38\,256\,613\,165&\(E_{\max}\)&1\,032\,928\,591\,041\\
\(H_C\)&1\,951\,087\,339\,951&\(T_1\)&4\,016\,944\,522\,473\\
\(T_2\)&8\,148\,658\,886\,639&\(b_{\min}\)&688\,618\,948\,863\\
\bottomrule
\end{tabular}
\end{center}
Thus
\begin{equation}\label{eq:smaller-d}
 b=5^{17},\qquad L_0=5^{19},\qquad
 d=5^{36}=14\,551\,915\,228\,366\,851\,806\,640\,625.
\end{equation}
The coefficient \(\beta\) has exactly \(d\) binary digits.
The corresponding package-06 value was \(5^{67}\); the reduction
factor is \(5^{31}\). This improves coefficient height but does not
make expansion feasible. The new recipe also removes four copy
coordinates and eight families of overlap tests.

\subsection{Populations, carry bounds, and exact convolution tests}

\begin{lemma}\label{lem:cross-interface}
The coefficients above satisfy
\begin{gather*}
 0<M_C,M_F<B-1,\qquad M_C\equiv2\pmod4,\qquad
 M_F\equiv4\pmod8,\\
 \pc(M_C)+\pc(M_F)=d.
\end{gather*}
They implement the four cross overlaps of
Section~\ref{sec:local} with the same carry-free native interface
as the package-06 compiler.
\end{lemma}

\begin{proof}
The powers of \(A_0\) occupy disjoint \(\ell\)-bit blocks, so
\[
 \pc(\mu)=\ell+5a_0+3+p_0,\qquad
 m_0=2\pc(\mu)+4a_0+2.
\]
The two clause sites in \(M_F^0\) contribute \(2\pc(\mu)\);
the four copy-test families contribute \(4a_0\); and the anchor
sites contribute two. Hence \(\pc(M_F^0)=m_0\) and
\(\pc(M_F)=m_0+1\). The added low bit \(4\) is disjoint from
all other support.

In \(M_C\), exactly \(n_0\) distinct bits are removed from \(B-1\):
one low bit at every native position except selector \(1\), plus
one upper bit at the first dummy. No borrow occurs. Thus
\[
 \pc(M_C)=d-m_0-1,\qquad
 \pc(M_C)+\pc(M_F)=d.
\]
Selector zero is permitted, selector one is forbidden, ordinary
native digits are Boolean, and the first dummy digit may be
\(0,1,2,3\). The lowest mask bits give the displayed congruences;
unused high positions give the strict range bounds.

The clause coefficients obey the conservative estimate
\[
 \sum_{e\in\mathcal E}c_e
       \le14(\kappa+5a_0+4)A_0^{5a_0+4},
 \qquad \mu<A_0^{j_{\max}+1}.
\]
Consequently the radix definition implies
\begin{equation}\label{eq:cross-carry}
 \rho\ge
 \max\left\{4(n_0+2)\left(2\sum_ec_e+6\right)+8,\
             2\mu+4,\ 16\right\}.
\end{equation}
This bounds every unshifted native field coefficient by
\(\rho/4-2\), including the permitted upper dummy bit.
Moreover \(T_2+E_{\max}<L_0\), so no unshifted field term
spills across a cell boundary.

At each clause site \(T_j\), convolution by \(D_C\) reads exactly
\(\sum_ec_ex_e\) from the current cell. The other clause band cannot
reach it because \(T_2-T_1>2E_{\max}\). The short-shift terms lie
strictly below the clause bands.

At \(e(0,1,u)\) and \(e(1,1,u)\), the short shift \(3a_0\)
reads, respectively, the current center and following copies.
The identity contribution from the temporal neighbor reads its
prior and center copies. These are the two temporal overlaps.

At \(H_C+e(1,0,u)\) and \(H_C+e(1,1,u)\), the shift \(H_C+a_0\)
reads the current center and right copies, whereas \(D_R\)
reads the spatial neighbor's left and center copies. These are
the two spatial overlaps.

At \(9M_0\) and \(27M_0\), the shifts \(8M_0\) and \(24M_0\)
read current anchors \(M_0\) and \(3M_0\); the temporal neighbor
contributes its anchors \(9M_0\) and \(27M_0\).
There are no unintended terms. Low support is below \(M_0\);
the four anchors are separated by powers of three; high clause
bands lie above the test sites; and
\(H_C=51M_0+3a_0+1\) separates horizontal tests from the possible
anchor sum \(24M_0+27M_0=51M_0\).
These support statements prove the exact convolution roles for
every word, not merely for the toy audits.
\end{proof}

\subsection{Why arbitrary bit rotations still synchronize}

We give the synchronization argument because it is the principal
compiler obligation beyond the population identity.
After arithmetic soundness recovers \(X=2^R\), write its binary
rotation exponent as \(bt+\epsilon\), \(0\le\epsilon<b\).
The possible rotated radix support is contained in a translate,
modulo \(L_0\), of
\[
 \mathcal E\cup\{e_*+1\}\subset[0,E_{\max}+1].
\]
The extra position records the upper dummy bit spilling into the
next radix position. Both cyclic distances between \(T_1\) and
\(T_2\) exceed \(E_{\max}+1\), so this rotated support cannot
contaminate both clause sites. One clause site is clean.

For \(\epsilon\le b-2\), a rotated digit is at most \(3\rho/4\).
For \(\epsilon=b-1\), it is at most \(\rho/2+1\).
Together with \eqref{eq:cross-carry}, neither case creates a carry.
Thus the mask tests apply to the actual field. More explicitly,
the transport equation and \(X=wq\equiv w\pmod{q-1}\) identify
\(F\) modulo \(q-1\) with the native field \(\widehat F\).
The carry bounds give \(0<\widehat F<q-1\), while arithmetic
soundness gives \(0<F<q\); therefore \(F=\widehat F\).

At the clean clause site, the low base-\(A_0\) digit counts
selected crosses. This count is at most \(\kappa<A_0-1\), and
its AND with \(A_0-2\) vanishes only for counts zero and one.
Every higher clause digit is a sum of the relevant selected
indicator and one copy or anchor bit. These sums are below \(A_0\),
so no base-\(A_0\) carry occurs. Their masked low bits force
evenness. Hence an occupied cell has exactly one selector, all
five matching copies, and all four anchors; an unoccupied cell
has none of these copies or anchors. Dummy bits have coefficient
zero and remain independent.

The marked origin is occupied. Its two anchor tests require
two distinct rotated digits to be odd. If
\(1\le\epsilon\le b-2\), all rotated digits are even.
If \(\epsilon=b-1\), only the single upper-dummy spill position
can be odd. Thus \(\epsilon=0\).

All nonanchor native positions lie below \(M_0\), and the only
pair of native positions separated by \(18M_0\) is
\((9M_0,27M_0)\). Since \(L_0>54M_0\), reduction modulo \(L_0\)
cannot supply another such pair through wraparound.
The two anchor tests therefore force \(t=0\pmod{L_0}\):
the surviving rotation is by a whole number of cells.

The copy tests now propagate occupancy to both spatial neighbors.
Indeed, a selected cross has one-hot copies, and its tested
center or side copy cannot match an empty neighbor. The spatial
offset \(1\) generates the whole cyclic word, so Start forces
every cell occupied. The four decoded overlaps then yield
\eqref{eq:reconstruct}. In particular a Start-containing solution
cannot have zero temporal stride, since its prior cell and center
are different. No uniqueness of Start, or independent End marker,
is required for soundness.

\subsection{Completeness at the actual packed index}

The preceding support argument also preserves the canonical
completeness construction of package 06~\cite{blank06}. We record
the index-control step explicitly, rather than assuming that a
formal congruent exponent can replace the actual packed value.

Choose a padded halting tableau with spatial width \(h_0\) and
temporal height \(H_0\), both sufficiently large powers of five,
such that \(H_0\ge25\) and \(N=h_0H_0>25d\).
The helical quotient has one Start at the origin. Set \(q=B^N\).
Initially set every dummy bit to zero and let \(C_0,F_0,R_0\)
be the actual native word, actual field, and actual packed index.
Switching on the low dummy at an eligible cell \(v\) changes
\begin{align}
 \delta C&=\rho^{e_*}B^v,\notag\\
 \delta F&=(D_C+BD_R+B^{h_0})\rho^{e_*}B^v,\notag\\
 \delta R&=-\Gamma B^v,\qquad
 \Gamma=\rho^{e_*}(q^2-1)
             [1+q(D_C+BD_R+B^{h_0})].
 \label{eq:cross-dummy-derivative}
\end{align}
Use distinct positions
\[
 \{4j:0\le j<N/5\}\cup\{1\}.
\]
Neither spatial nor temporal movement wraps these individual
contributions, because \(h_0=N/H_0\le N/25\).

The two clause-band terms in \(D_C\) cancel modulo five:
\[
 T_2-T_1\equiv2\pmod4,\qquad \rho\equiv2\pmod5.
\]
Since \(a_0\equiv3\pmod4\),
\begin{equation}\label{eq:cross-unit}
 2D_C-D_R
 \equiv2(2+3\rho^{H_C}+1+1)-\rho^{H_C}
 \equiv3\pmod5.
\end{equation}
Also \(q,B,B^{h_0}\equiv2\pmod5\), so the bracket in
\eqref{eq:cross-dummy-derivative} is \(2D_C-D_R\) modulo five.
Thus \(\Gamma\) is invertible modulo the power of five \(dN\).
No conditional correction monomial is needed.

Every residue modulo \(dN\) is a Boolean subset sum of the displayed
weights \(B^v\). To see this, \(B^4\) has order \(N/5\) modulo
\(dN\), by
\[
 \nu_5(2^{4dj}-1)=\nu_5(d)+1+\nu_5(j).
\]
Its powers enumerate exactly the residues \(1+5dt\), with
\(t\) modulo \(N/5\). For a desired residue, optionally subtract
the distinct weight \(B\) so that the remaining residue is
nonzero modulo five. Let \(k_0\) be its least positive residue
modulo \(5d\). Then \(1\le k_0<5d<N/5\) and \(k_0\) is a unit
modulo \(N/5\). A cyclic interval of \(k_0\) distinct residues \(t\)
has sum \(k_0t_0+k_0(k_0-1)/2\); varying \(t_0\) gives every
sum modulo \(N/5\). Choosing the corresponding powers of \(B^4\)
proves the subset assertion.

Apply it to
\((R_0-dh_0)\Gamma^{-1}\pmod{dN}\).
After this Boolean adjustment the \emph{actual} packed index obeys
\begin{equation}\label{eq:cross-tuned-index}
 R\equiv dh_0\pmod{dN}.
\end{equation}
The dummy is absent from every clause and copy test; its upper bit
remains zero. All masks survive the adjustment. The canonical
digit bounds also give
\[
 2C+F\le(\rho/4+1)\frac{q-1}{\rho-1}<q/3,
\]
so the outer slack \(q-F-C\) stays strictly positive.

The shifted field mask deserves a separate check. Its unit radix
digit comes only from the temporal neighbor's Start selector,
since \(D_C,D_R\) have positive minimum exponent. Canonical Start
occurs only at zero and \(0<h_0<N\). That digit is therefore zero
at the origin and zero or one elsewhere. The origin's two low
mask bits and the other cells' bit \(2\) vanish as required.
Dummy adjustment occurs at positive radix exponent and changes
none of these tests.

Consequently the inverse-population identity and the native masks
give, at this same final \(R\),
\[
 3q+1<R<q^4,\qquad
 R\equiv3\pmod4,\qquad
 \pc(R)=3dN+2.
\]
The positive half-binomial converse~\cite{halfbinomial,blank06}
now supplies fresh Pell witnesses at the actual index.
It gives \(X=2^R\), \(q^3\mid Y\), and hence \(w=X/q\in\Pos\).
Equation \eqref{eq:cross-tuned-index} yields
\(w\equiv B^{h_0}\pmod{q-1}\). The native transport quotient
\[
 z=\frac{(K+w)C-F}{q-1}
\]
is integral and positive, since \(w>q>F\) and \(C\ge1\).
This is the shifted transport convention of
\eqref{eq:system60}. The ordinary strong witness is obtained
from the canonical normalized one by multiplication by \(\Delta\).
Thus the smaller coefficient quartet preserves completeness.

\subsection{Scope and reproducible evidence}

The inherited dependencies are the marked helical-tableau argument,
the native carry and mask interpretation, and the half-binomial
Pell converse~\cite{blank06,raw75,halfbinomial}. The archive also
preserves the native 76-operation compiler proof, the half-binomial
mask correction, and the five-adic subset-control note, with exact
source identities. The present change is the five-coordinate
language and its revised finite layout; it does not silently
specialize a theorem whose input-marker hypotheses remain unmet.

The generator \code{compiler/cross_constants.py} emits the closed
recipe, all scalar geometry, and its audits. With \code{--tm}, it
checks the exact source-table hash, all thirteen census categories,
32 helical control presentations, and 1,144 projected crosses.
Its other checks include 26,106 cyclic lift checks, 100,384 exhaustive
small cross-word candidates, and 16,224 sparse convolution basis
tests. An independent audit rederives the large-instance geometry
and support isolation by interval arithmetic without importing the
generator.

These finite checks supplement the parametric proofs above.
They do not enumerate the actual selector alphabet, expand the
coefficients, exhibit a halting K215 run, or kernel-check the
Diophantine equivalence. All definitions and proofs are ordinary
finite mathematics formalizable in ZFC; supplying such an argument
is distinct from supplying an encoded ZFC derivation or a new
Lean/Rocq theorem.

